% !PRINTED 673
\weberpage{673}
This simple remark yields the following important theorem:

2. \gesperrt{In an algebraic field $\Omega$ there is only a finite number of integers $\eta$ of such a nature that the absolute values of the conjugate numbers $\eta_s$ lie below a given limit.}

We now make use of the general theorem \S.\,155, (25), according to which, if $f(x_1,x_2,\dots,x_n)$ is any \gesperrt{positive quadratic form} in $n$ variables with the determinant $D$, the variables $x_i$ can be determined as integral rational numbers, not all vanishing, in such a way that
\[
f(x_1,x_2,\dots,x_n) < n\sqrt[n]{0{,}404\,D},
\]
i.\,e.\ so that it becomes smaller than a number depending on the determinant $D$ alone.

We apply this theorem, as in \S.\,156, to the form:
\[
f(x_1,x_2,\dots,x_n) = \frac{|\eta_1|^2}{c_1^2} + \frac{|\eta_2|^2}{c_2^2} + \cdots + \frac{|\eta_n|^2}{c_n^2},
\]
in which the $\eta_s$ are the linear forms (4), and, in the case of a real $\eta_s$,
\[
|\eta_s|^2 = (\eta_s)^2,
\]
while in the case of an imaginary pair $\eta_s, \eta_s'$
\[
|\eta_s|^2 = |\eta_s'|^2 = \eta_s\eta_s'.
\]
In the latter case $\eta_s\eta_s'$ is the sum of two real squares, and accordingly, if the $c_s$ are assumed to be real, $f$ is a positive form. Concerning the real quantities $c_s$, which have so far been arbitrary, we further stipulate that, if $\eta_s, \eta_s'$ form a conjugate imaginary pair,
\weq{5}{c_s = c_s'}
shall hold; besides this we shall subject the $c_s$ to the further condition:
\weq{6}{c_1 c_2 \dots c_n = 1.}
We form the determinant of the function $f$, as in \S.\,156, by first regarding it as a form in the variables $\eta_1, \eta_2, \dots, \eta_n$, and the determinant of this form is, on account of (6), equal to $\pm 1$. The determinant of $f$ with respect to the variables $x_i$ is therefore (Bd.\ I, \S.\,56), apart from the sign, equal to the
%CONT>>
